\scanpage{93}
\mspage{99}

Dans certaines questions il suffira de savoir que ledit foncteur est pleinement fidèle. Cela \deleted{équivaut sans doute}\inserted{implique} \deleted{à dire} que

\[
H^{\unsure{\cdot|i}}_{\mathbb D}(Z,F)\overset{\sim}{\unsure{\relbar|\longrightarrow}}H^{\cdot}_{\mathbb D}(Z',f^*F)
\]

quand $F$ est \uline{loc.} \uline{constant} sur $Z$ \deleted{et} \deleted{\unsure{(quelle)}} (mais il n'est pas clair que cette condition implique tjrs la pleine fidélité), de $=$ pour l'homologie

\[
H^{\mathbb D}_{\cdot}(Z',\unsure{f'|f^*}F)\longrightarrow H^{\mathbb D}_{\cdot}(Z,F)
\]

% slanted
\margin{Mais c'est OK si on a des $\mathrm{Hom}$ internes dans les $\mathbb D(Z)\ \text{loc. const.}$, satisfaisant $\mathrm{Hom}(F,G)\simeq\mathrm{Hom}_{\unsure{C|e}}(e_Z,\mathrm{Hom}(F,G))$, où $e$ est l'élément final de $A$, les $\mathrm{Hom}$ devant \unsure{compatibles} aux images inverses.}

% slanted
\margin{On \unsure{pourra|pourrait} dire que $f:Z'\to Z$ est \deleted{alors} $\mathbb D$-équiv. \uline{0-complète} \uline{à gche}}

% slanted
\margin{et \deleted{\unsure{à droite}}\inserted{0-complète} \uline{à droite} dans le cas homologique.}

Moyennant cette condition, on a la

\uline{Proposition} : \inserted{On suppose \uline{Der} 3 b' satisfait} Soit $f:X\to Y$ dans $\mathrm{Diag}$. Pour que $f$ soit une $\mathbb D$-fibration, il f. et s. que pour $Y''\xrightarrow{p}Y'$ sur $Y$ qui soit une $\mathrm{Hot}$-équivalence, $X''=X\times_Y Y''\unsure{\relbar|\to} X'=X\times_Y Y'$ soit \deleted{une} $\mathbb D$-équivalence.

À vrai dire, on peut énoncer quelque chose de raisonnable, sans supposer l'axiome Der 3 b'. On doit introduire la

\uline{Définition} : Soit $f:Z'\unsure{\relbar|\to}Z$ dans $\mathrm{Diag}$. On dit que $f$ est une $\mathbb D$-équivalence 1-\uline{complète} si le foncteur $f^*:\mathbb D(Z)\unsure{\relbar|\to}\mathbb D(Z')$ est pleinement fidèle sur la sous-catégorie pleine de $\mathbb D(Z)$ formée des coeff\textsuperscript{ts} \uuline{loc.} constants. On dit que $f$ est une $\mathbb D$-équivalence 2-\uline{complète} si elle induit une équiv. entre les cat. de coeff\textsuperscript{ts} loc. constants.

% archive sheet number: 92
